% Chapter 3: System Design
\chapter{System Design}
\label{ch:design}

This chapter presents the architectural decisions, data model, security design,
and external-service integrations that constitute OrderVista.  Each design
choice is discussed alongside the alternatives considered.

% ─────────────────────────────────────────────────────────────
\section{High-Level Architecture}
\label{sec:architecture}

OrderVista follows a four-layer client--server architecture
(Figure~\ref{fig:arch-layers}):

\begin{enumerate}
  \item \textbf{Client layer} — a single Flutter codebase compiled to Android,
        containing both the admin portal and driver application as role-gated
        views.
  \item \textbf{Backend-as-a-Service (BaaS) layer} — Supabase, providing
        PostgreSQL, authentication, real-time subscriptions, and Supabase
        Storage for media.
  \item \textbf{External services layer} — Mapbox for geocoding and static map
        tiles; OpenRouter for proxying DeepSeek~R1 inference requests.
  \item \textbf{Security layer} — cutting across all tiers via JWT-based
        authentication, PostgreSQL Row-Level Security (RLS) policies, and
        role-based access control (RBAC).
\end{enumerate}

\begin{figure}[ht]
  \centering
  \resizebox{0.92\textwidth}{!}{%
  \begin{tikzpicture}[
    layer/.style={draw, rounded corners, minimum width=11cm, minimum height=1.2cm,
                  font=\sffamily\small, text=white, align=center},
    proto/.style={font=\sffamily\footnotesize, text=black},
    arr/.style={-{Stealth[length=3mm]}, thick, gray!70},
    node distance=1.1cm
  ]
    % Layers
    \node[layer, fill=teal!80!black]                       (client)   {Client Layer \\ \footnotesize Flutter 3.35 / Android — Admin Portal \& Driver App};
    \node[layer, fill=blue!70!black, below=of client]      (baas)     {Backend-as-a-Service Layer \\ \footnotesize Supabase: Auth, Realtime (WebSocket), Storage, PostgREST};
    \node[layer, fill=violet!70!black, below=of baas]      (db)       {Database Layer \\ \footnotesize PostgreSQL 15 + PostGIS + Row-Level Security};
    \node[layer, fill=orange!80!black, below=of db]        (ext)      {External Services \\ \footnotesize Mapbox Geocoding \quad|\quad OpenRouter / DeepSeek R1};

    % Security overlay
    \node[draw=red!60, dashed, very thick, rounded corners,
          fit=(client)(baas)(db)(ext),
          inner sep=8pt, label={[red!60, font=\sffamily\footnotesize]right:Security Layer (JWT + RLS + RBAC)}] {};

    % Protocol arrows
    \draw[arr] (client) -- node[proto, right, xshift=2pt] {HTTPS + WSS} (baas);
    \draw[arr] (baas)   -- node[proto, right, xshift=2pt] {SQL (PostgREST)} (db);
    \draw[arr] (db)     -- node[proto, right, xshift=2pt] {REST API} (ext);
  \end{tikzpicture}}%
  \caption{Four-layer architecture of OrderVista with cross-cutting security layer.}
  \label{fig:arch-layers}
\end{figure}

\paragraph{Alternatives considered.}
A traditional three-tier stack (e.g.\ Express.js + PostgreSQL + React Native)
was evaluated first.  Supabase was chosen instead because its hosted
PostgreSQL already includes Auth, Realtime, and Storage, eliminating the need
to write and deploy a custom API server.  This reduced the project scope to a
single-developer effort while preserving enterprise-grade security through
native RLS support \autocite{supabase2026docs}.  Firebase was rejected because
its document-oriented Firestore model lacks native RLS and makes relational
queries (e.g.\ joining routes to tasks to drivers) cumbersome
\autocite{abdul2018multitenant}.

\subsection{Deployment Architecture}
\label{sec:deployment}

Figure~\ref{fig:deployment} shows the physical deployment topology.  All
server-side components are managed services, eliminating operational overhead
for infrastructure provisioning.

\begin{figure}[ht]
  \centering
  \resizebox{0.94\textwidth}{!}{%
  \begin{tikzpicture}[
    node/.style={draw, rounded corners=4pt, minimum width=2.8cm,
                 minimum height=0.8cm, font=\sffamily\footnotesize,
                 align=center, thick},
    device/.style={node, fill=teal!15},
    cloud/.style={node, fill=blue!10},
    ext/.style={node, fill=orange!10},
    arr/.style={-{Stealth[length=2.5mm]}, thick, gray!60},
    darr/.style={{Stealth[length=2.5mm]}-{Stealth[length=2.5mm]}, thick, gray!60},
    lbl/.style={font=\sffamily\tiny, text=gray!80},
    node distance=1.2cm and 2.4cm
  ]
    % Client devices
    \node[device]                            (android) {Android Device\\{\tiny Samsung Galaxy A54}};
    \node[device, right=3.5cm of android]    (web)     {Chrome Browser\\{\tiny Admin Portal}};

    % Supabase cloud
    \node[cloud, below=2cm of android]       (auth)    {Supabase Auth\\{\tiny JWT issuance}};
    \node[cloud, right=of auth]              (rest)    {PostgREST\\{\tiny REST $\to$ SQL}};
    \node[cloud, right=of rest]              (rt)      {Realtime\\{\tiny WebSocket}};
    \node[cloud, right=of rt]               (store)   {Supabase Storage\\{\tiny PoD images}};

    % Database
    \node[cloud, below=1.5cm of rest, fill=blue!18] (pg) {PostgreSQL 15\\{\tiny RLS + PostGIS}};

    % External
    \node[ext, left=of pg]                   (mapbox)  {Mapbox API\\{\tiny Geocoding, Static Maps}};
    \node[ext, right=of pg]                  (or)      {OpenRouter API\\{\tiny DeepSeek R1 / GPT-4o-mini}};

    % Local DB
    \node[device, below=1cm of android, xshift=-1cm]  (drift)  {Drift SQLite\\{\tiny Local cache}};

    % Grouping
    \begin{scope}[on background layer]
      \node[draw=blue!30, dashed, rounded corners, fill=blue!3,
            fit=(auth)(rest)(rt)(store)(pg),
            inner sep=10pt, label={[blue!50, font=\sffamily\footnotesize]above:Supabase Cloud}] {};
    \end{scope}

    % Connections
    \draw[darr] (android) -- node[lbl, left] {HTTPS} (auth);
    \draw[darr] (web) -- node[lbl, right] {HTTPS} (rt);
    \draw[arr]  (android) -- (drift);
    \draw[arr]  (auth) -- (pg);
    \draw[arr]  (rest) -- (pg);
    \draw[arr]  (rt) -- (pg);
    \draw[arr]  (store) -- (pg);
    \draw[darr] (pg) -- (mapbox);
    \draw[darr] (pg) -- (or);
  \end{tikzpicture}}%
  \caption{UML deployment diagram showing physical nodes and communication protocols.}
  \label{fig:deployment}
\end{figure}

% ─────────────────────────────────────────────────────────────
\section{Technology Stack}
\label{sec:tech-stack}

Table~\ref{tab:tech-stack} summarises the core technologies and the rationale
for each selection.

\begin{table}[ht]
  \centering
  \caption{Technology stack with alternatives considered.}
  \label{tab:tech-stack}
  \small
  \begin{tabularx}{\textwidth}{l l X}
    \toprule
    \textbf{Layer} & \textbf{Choice} & \textbf{Rationale / Alternatives} \\
    \midrule
    Mobile UI
      & Flutter 3.35
      & Single codebase for Android \& iOS; Skia renderer ensures pixel-perfect
        maps.  React Native was considered but its bridge-based architecture
        adds latency for map-heavy views
        \autocite{flutter2026docs}. \\
    Backend
      & Supabase (hosted)
      & Built-in Auth, Realtime, Storage, and PostgREST.  A custom Node.js API
        would require separate deployment and maintenance
        \autocite{supabase2026docs}. \\
    Database
      & PostgreSQL 15
      & Native RLS, PostGIS for spatial queries, JSONB for flexible telemetry
        payloads.  MySQL lacks RLS; MongoDB lacks relational joins
        \autocite{postgresql2026rls}. \\
    Geocoding
      & Mapbox Geocoding API
      & 100\,000 free requests/month; returns GeoJSON with structured address
        components.  Google Geocoding API was costlier for the same tier
        \autocite{mapbox2026geocoding}. \\
    AI / LLM
      & DeepSeek R1 / GPT-4o-mini via OpenRouter
      & Configurable model backend; DeepSeek~R1 provides chain-of-thought
        reasoning at \$0.14/M input tokens; GPT-4o-mini serves as a
        low-latency fallback.  Both accessed through OpenRouter's unified
        API gateway \autocite{deepseek2025r1,openrouter2026docs}. \\
    Local DB
      & Drift (SQLite)
      & Type-safe, reactive, compiles to native SQLite; enables offline-first
        with sync queue.  Hive lacks SQL expressiveness
        \autocite{drift2026docs}. \\
    Version Control
      & Git + GitHub
      & Branch-per-feature workflow; CI via GitHub Actions for APK builds. \\
    \bottomrule
  \end{tabularx}
\end{table}

% ─────────────────────────────────────────────────────────────
\section{Data Model}
\label{sec:data-model}

The PostgreSQL schema consists of thirteen tables distributed across three migrations.
Table~\ref{tab:schema} lists each table and its purpose; foreign-key
relationships are described in the text below.

\begin{table}[ht]
  \centering
  \caption{Database schema overview (13 tables).}
  \label{tab:schema}
  \small
  \begin{tabularx}{\textwidth}{l l X}
    \toprule
    \textbf{Table} & \textbf{PK Type} & \textbf{Purpose} \\
    \midrule
    \texttt{companies}          & UUID & Multi-tenant root; stores name and timezone. \\
    \texttt{company\_members}   & UUID & Maps \texttt{auth.users} to companies with a role
                                         (\texttt{driver}, \texttt{dispatcher}, \texttt{admin}). \\
    \texttt{vehicles}           & UUID & Fleet registry; links to an active route and assigned driver. \\
    \texttt{routes}             & UUID & Route header; references company, vehicle, driver; tracks
                                         status, optimizer type, and ETA summary (JSONB). \\
    \texttt{tasks}              & UUID & Individual delivery tasks with address, PostGIS
                                         \texttt{geography(point,4326)} location, priority, and
                                         proof-of-delivery fields. \\
    \texttt{route\_tasks}       & UUID & Join table; assigns a task to a route with a sequence
                                         number and planned arrival/departure. \\
    \texttt{telemetry\_events}  & BIGINT & GPS pings from drivers; includes speed, heading,
                                            battery, and a JSONB payload for extensibility. \\
    \texttt{sync\_queue\_items} & UUID & Offline-first write queue; stores pending mutations with
                                          retry count and backoff timestamp. \\
    \texttt{media\_artifacts}   & UUID & Photos and documents (PoD images); references Storage
                                          paths with optional checksums and expiry. \\
    \texttt{chat\_messages}     & UUID & Admin--driver messaging; supports directed and broadcast
                                          messages with read-receipt tracking. \\
    \texttt{notifications}      & UUID & System notifications for task assignments, status changes,
                                          and admin broadcasts. \\
    \texttt{geofence\_rules}    & UUID & Configurable geofence boundaries for delivery zones;
                                          supports future proximity-based alerts. \\
    \texttt{audit\_log}         & UUID & Immutable log of security-relevant events (logins, role
                                          changes, RLS denials) for compliance auditing. \\
    \bottomrule
  \end{tabularx}
\end{table}

\paragraph{Key relationships.}
Every table carries a \texttt{company\_id} foreign key, enforcing
tenant-scoped partitioning at the schema level.  \texttt{routes} reference
both a \texttt{vehicle} and a \texttt{driver} (both nullable for draft
routes).  \texttt{route\_tasks} forms the many-to-many join between routes
and tasks, with a unique compound constraint on
\texttt{(route\_id, sequence)} to guarantee ordering.  \texttt{tasks} stores
the PostGIS \texttt{geography} column used by the Haversine-based
nearest-neighbour sequencer at route-generation time.
Figure~\ref{fig:er} illustrates the principal foreign-key relationships.

\begin{figure}[ht]
  \centering
  \resizebox{0.92\textwidth}{!}{%
  \begin{tikzpicture}[
    entity/.style={draw, rectangle, rounded corners=2pt,
                   minimum width=2.4cm, minimum height=0.7cm,
                   font=\sffamily\footnotesize, fill=blue!8},
    rel/.style={-{Stealth[length=2.5mm]}, thick, gray!60},
    node distance=1.6cm and 2.2cm
  ]
    % Core entities
    \node[entity]                            (comp)  {companies};
    \node[entity, below left=of comp]        (mem)   {company\_members};
    \node[entity, below right=of comp]       (veh)   {vehicles};
    \node[entity, below=2.4cm of comp]       (route) {routes};
    \node[entity, below left=of route]       (task)  {tasks};
    \node[entity, below=of route]            (rtask) {route\_tasks};
    \node[entity, right=of route]            (telem) {telemetry\_events};
    \node[entity, left=of task]              (media) {media\_artifacts};
    \node[entity, right=of telem]            (sync)  {sync\_queue};
    \node[entity, below=of media]            (chat)  {chat\_messages};
    \node[entity, below=of sync]             (notif) {notifications};
    \node[entity, below right=of chat]       (audit) {audit\_log};
    \node[entity, below=of task]             (geo)   {geofence\_rules};

    % Foreign keys
    \draw[rel] (mem)   -- (comp);
    \draw[rel] (veh)   -- (comp);
    \draw[rel] (route) -- (comp);
    \draw[rel] (route) -- (veh);
    \draw[rel] (route.west) -- (mem);
    \draw[rel] (task)  -- (comp);
    \draw[rel] (rtask) -- (route);
    \draw[rel] (rtask) -- (task);
    \draw[rel] (telem) -- (mem);
    \draw[rel] (media) -- (task);
    \draw[rel] (sync)  -- (mem);
    \draw[rel] (chat)  -- (mem);
    \draw[rel] (notif) -- (comp);
    \draw[rel] (audit) -- (comp);
    \draw[rel] (geo)   -- (comp);
  \end{tikzpicture}}%
  \caption{Entity-relationship overview — arrows denote foreign-key references.}
  \label{fig:er}
\end{figure}

\paragraph{Alternatives considered.}
A NoSQL document model (e.g.\ Firestore collections) would have simplified
schema evolution but would have required application-level joins and forfeited
PostGIS spatial indexing.  The relational model was preferred because the core
domain—routes comprising ordered tasks assigned to drivers—is inherently
relational, and PostGIS enables efficient distance calculations directly in
SQL \autocite{postgresql2026rls}.

% ─────────────────────────────────────────────────────────────
\section{Security Design}
\label{sec:security}

Security in OrderVista is enforced at three complementary levels:
authentication, row-level security, and role-based access control.

\subsection{Authentication}
\label{sec:authn}

Supabase Auth issues JWTs upon email/password login.  Every subsequent
PostgREST and Realtime request carries the JWT, which PostgreSQL validates
internally via the \texttt{auth.uid()} function.  This removes the need for
a custom authentication middleware \autocite{supabase2026docs}.

\subsection{Row-Level Security (RLS)}
\label{sec:rls}

Eleven of the thirteen tables have RLS enabled; \texttt{audit\_log}
(immutable, insert-only) and \texttt{geofence\_rules} (reserved for future
use) are exempt.  A shared helper function,
\texttt{is\_company\_member(company\_id)}, returns \texttt{true} only when the
authenticated user belongs to the specified company.  Every \texttt{SELECT},
\texttt{INSERT}, \texttt{UPDATE}, and \texttt{DELETE} policy on the eleven
RLS-protected tables invokes this function, guaranteeing that a query can
never read or modify rows belonging to a different tenant
\autocite{postgresql2026rls,aws2020rls}.

For chat messages, additional fine-grained policies restrict reads to the
sender, recipient, or broadcast scope, and restrict \texttt{UPDATE} (marking
messages as read) to the designated recipient.

\subsection{Role-Based Access Control (RBAC)}
\label{sec:rbac}

Table~\ref{tab:rbac} summarises the permission matrix.

\begin{table}[ht]
  \centering
  \caption{RBAC permission matrix (admin vs.\ driver).}
  \label{tab:rbac}
  \small
  \begin{tabularx}{\textwidth}{X c c}
    \toprule
    \textbf{Capability} & \textbf{Admin} & \textbf{Driver} \\
    \midrule
    Create / edit / delete tasks            & \checkmark & --- \\
    Import tasks via CSV                    & \checkmark & --- \\
    Generate and assign routes              & \checkmark & --- \\
    View all drivers and their locations    & \checkmark & --- \\
    Review proof-of-delivery submissions    & \checkmark & --- \\
    View assigned task list                 & \checkmark & \checkmark \\
    Update task status (in-progress, done)  & --- & \checkmark \\
    Capture PoD photo and signature         & --- & \checkmark \\
    Use AI assistant (summary \& chat)      & \checkmark & \checkmark \\
    Send / receive chat messages            & \checkmark & \checkmark \\
    \bottomrule
  \end{tabularx}
\end{table}

\paragraph{Alternatives considered.}
Application-level authorisation (e.g.\ middleware guards) was considered but
rejected because it cannot protect against direct database access or
misconfigured API endpoints.  Database-level RLS ensures that even if the API
layer is bypassed, data isolation is maintained—a well-established pattern
for multi-tenant SaaS systems \autocite{gao2011nonintrusive}.

\subsection{Task Lifecycle}
\label{sec:task-lifecycle}

Each delivery task follows a three-state lifecycle enforced by both the UI
and the database layer.  Figure~\ref{fig:task-fsm} visualises the permitted
transitions together with the actor responsible for each.

\begin{figure}[ht]
  \centering
  \resizebox{0.88\textwidth}{!}{%
  \begin{tikzpicture}[
    state/.style={draw, rounded corners=6pt, minimum width=2.6cm,
                  minimum height=1cm, font=\sffamily\small,
                  fill=blue!10, align=center, thick},
    init/.style={circle, fill=black, inner sep=2pt},
    arr/.style={-{Stealth[length=3mm]}, thick, gray!70},
    lbl/.style={font=\sffamily\footnotesize, text=black, fill=white,
                inner sep=2pt, rounded corners=2pt},
    node distance=3.8cm
  ]
    % States
    \node[init]                      (dot)       {};
    \node[state, right=2cm of dot]   (pending)   {\textbf{Pending}\\{\footnotesize Created by admin}};
    \node[state, right=of pending]   (transit)   {\textbf{In-Transit}\\{\footnotesize Driver en route}};
    \node[state, right=of transit]   (delivered) {\textbf{Delivered}\\{\footnotesize PoD captured}};

    % Transitions
    \draw[arr] (dot)       -- (pending);
    \draw[arr] (pending)   -- node[lbl, above] {Driver starts task} (transit);
    \draw[arr] (transit)   -- node[lbl, above] {Driver submits PoD} (delivered);

    % Re-assignment loop
    \draw[arr] (transit) to[out=-150,in=-30]
      node[lbl, below] {Admin reassigns} (pending);
  \end{tikzpicture}}%
  \caption{Task lifecycle state machine with permitted transitions.}
  \label{fig:task-fsm}
\end{figure}

% ─────────────────────────────────────────────────────────────
\section{External Service Integration}
\label{sec:external}

\subsection{Mapbox Geocoding and Static Maps}
\label{sec:mapbox}

When an admin creates a task, the address string is sent to the Mapbox
Geocoding API, which returns latitude/longitude coordinates stored as a
PostGIS \texttt{geography} point.  The geocoded location is subsequently
used by the nearest-neighbour route sequencer and by the static-map tile
generator, which renders a preview image of the planned route for the admin
dashboard \autocite{mapbox2026geocoding}.

\subsection{DeepSeek~R1 via OpenRouter}
\label{sec:deepseek}

The AI assistant calls DeepSeek~R1 through the OpenRouter proxy, which
provides a single API key and unified billing across multiple LLM providers.
Two prompt templates are employed:

\begin{itemize}
  \item \textbf{Summary prompt} — injects the current driver's task list,
        route status, and fleet statistics into a system prompt, requesting a
        concise natural-language summary.
  \item \textbf{Chat prompt} — extends the summary context with the user's
        free-text question and a JSON action schema, enabling the model to
        return structured actions (e.g.\ \texttt{\{"action": "re\_sequence",
        "params": \{"route\_id": "..."\}\}}) that the client can execute with
        a single confirmation tap.
\end{itemize}

This agentic pattern—where the LLM proposes actions and the user confirms—was
chosen over full autonomous execution to prevent hallucination-induced errors,
a known LLM reliability concern \autocite{ji2024hallucination}.

% ─────────────────────────────────────────────────────────────
\section{Critical System Flows}
\label{sec:flows}

\subsection{Administrative Workflow}
\label{sec:admin-flow}

The admin follows an eight-step workflow:

\begin{enumerate}
  \item \textbf{Task creation} — manual entry or bulk CSV import with
        validation (required fields: address, optional: time window, priority,
        notes).
  \item \textbf{Geocoding} — each task address is geocoded via Mapbox; failures
        are flagged for manual correction.
  \item \textbf{Route generation} — the nearest-neighbour heuristic sequences
        tasks by Haversine distance from a configurable depot.
  \item \textbf{Driver recommendation} — a composite score combining task
        proximity, current workload, historical on-time rate, and average
        completion time ranks available drivers.
  \item \textbf{Route assignment} — the admin reviews the recommended driver,
        map preview, and estimated distances, then confirms assignment.
  \item \textbf{Monitoring} — real-time Supabase subscriptions push driver
        location, task status changes, and telemetry events to the admin
        dashboard.
  \item \textbf{Mid-shift intervention} — if conditions change, the admin (or
        AI assistant) can trigger re-sequencing from the driver's current GPS
        position.
  \item \textbf{PoD review} — upon delivery completion, the admin reviews the
        uploaded photo, signature, and confirmation code.
\end{enumerate}

\subsection{Driver Workflow}
\label{sec:driver-flow}

\begin{enumerate}
  \item \textbf{Authentication} — email/password login; JWT issued.
  \item \textbf{Task list} — assigned tasks displayed in sequence order with
        real-time status badges.
  \item \textbf{Navigation} — tapping a task opens the address in the device's
        default maps application.
  \item \textbf{Status progression} — the driver marks each task as
        \emph{in-progress} and then \emph{completed}.
  \item \textbf{PoD capture} — mandatory photo and optional signature before
        a task can be marked complete.
  \item \textbf{Sync} — completed tasks are written to the local Drift
        database and enqueued for upload; the sync queue retries with
        exponential backoff if connectivity is lost.
  \item \textbf{AI assistant} — the driver can request a summary of remaining
        work or ask free-form questions about their route.
\end{enumerate}

\subsection{Offline Synchronisation}
\label{sec:offline-sync}

OrderVista follows an offline-first pattern.  All writes are persisted to the
local Drift (SQLite) database before being enqueued in the
\texttt{sync\_queue\_items} table.  A background isolate processes the queue
with a progressive backoff schedule (0, 10, 30, 60, 120\,s delays, up to five
attempts).  On reconnection, Supabase
Realtime subscriptions automatically push any server-side changes that
occurred while the device was offline, and the client reconciles state via
\texttt{updated\_at} timestamps.  Offline-first design is widely recommended
for mobile applications that must remain usable under unreliable network
conditions; modern frameworks commonly employ local caching and deferred
synchronisation to provide an ``always-on'' user experience
\autocite{gaur2024offlinefirst}.  This design ensures that drivers operating
in areas with intermittent cellular coverage, which is common in Hong Kong's
industrial estates, can continue capturing deliveries without data loss.

% ─────────────────────────────────────────────────────────────
\section{Interaction Sequences}
\label{sec:sequences}

This section formalises the two most architecturally significant interactions
as UML-style sequence diagrams: task creation with route generation
(Figure~\ref{fig:seq-task}) and the agentic AI suggestion cycle
(Figure~\ref{fig:seq-agentic}).

\subsection{Task Creation and Route Assignment}
\label{sec:seq-task}

Figure~\ref{fig:seq-task} traces the complete lifecycle of a task batch from
CSV import through route generation to driver notification.

\begin{figure}[ht]
  \centering
  \resizebox{\textwidth}{!}{%
  \begin{tikzpicture}[
    actor/.style={draw, rectangle, rounded corners=3pt, minimum width=2cm,
                  minimum height=0.7cm, font=\sffamily\footnotesize,
                  fill=#1!15, align=center, thick},
    msg/.style={-{Stealth[length=2mm]}, thick, font=\sffamily\tiny},
    ret/.style={-{Stealth[length=2mm]}, thick, dashed, font=\sffamily\tiny},
    note/.style={draw, rectangle, rounded corners=2pt, fill=yellow!10,
                 font=\sffamily\tiny, align=left, inner sep=3pt},
  ]
    % Participants
    \node[actor=teal]   (admin)  at (0,0)    {Admin\\Browser};
    \node[actor=orange]  (flutter) at (3.2,0)  {Flutter\\AdminPortal};
    \node[actor=violet]  (mapbox)  at (6.4,0)  {Mapbox\\API};
    \node[actor=blue]   (supa)   at (9.6,0)  {Supabase\\PostgREST};
    \node[actor=blue]   (pg)     at (12.8,0) {PostgreSQL\\(RLS)};
    \node[actor=teal]   (driver) at (16,0)   {Driver\\App};

    % Lifelines
    \foreach \n in {admin, flutter, mapbox, supa, pg, driver} {
      \draw[dashed, gray!40] (\n.south) -- ++(0,-12.5);
    }

    % Messages
    \def\y{-1.2}
    \draw[msg] (admin) -- node[above] {Upload CSV file} (flutter |- 0,\y);
    \def\y{-2.0}
    \draw[msg] (flutter |- 0,\y) -- node[above] {Parse \& validate rows} (flutter |- 0,\y);
    \def\y{-2.8}
    \draw[msg] (flutter |- 0,\y) -- node[above] {geocode(address)} (mapbox |- 0,\y);
    \def\y{-3.6}
    \draw[ret] (mapbox |- 0,\y) -- node[above] {lat, lng} (flutter |- 0,\y);
    \def\y{-4.4}
    \draw[msg] (flutter |- 0,\y) -- node[above] {INSERT tasks (batch)} (supa |- 0,\y);
    \def\y{-5.0}
    \draw[msg] (supa |- 0,\y) -- node[above] {RLS policy check} (pg |- 0,\y);
    \def\y{-5.6}
    \draw[ret] (pg |- 0,\y) -- node[above] {OK (rows inserted)} (supa |- 0,\y);
    \def\y{-6.2}
    \draw[ret] (supa |- 0,\y) -- node[above] {201 Created} (flutter |- 0,\y);

    % Route generation
    \def\y{-7.2}
    \node[note] at (-1.5, \y) {Admin clicks\\``Generate Route''};
    \draw[msg] (admin |- 0,\y) -- node[above] {generateRoute(taskIds)} (flutter |- 0,\y);
    \def\y{-8.0}
    \draw[msg] (flutter |- 0,\y) -- node[above, align=center] {nearestNeighbour\\Sequence(tasks)} (flutter |- 0,\y);
    \def\y{-8.8}
    \draw[msg] (flutter |- 0,\y) -- node[above] {INSERT route + route\_tasks} (supa |- 0,\y);
    \def\y{-9.4}
    \draw[msg] (supa |- 0,\y) -- node[above] {RLS check} (pg |- 0,\y);
    \def\y{-10.0}
    \draw[ret] (pg |- 0,\y) -- node[above] {OK} (supa |- 0,\y);
    \def\y{-10.6}
    \draw[ret] (supa |- 0,\y) -- node[above] {route created} (flutter |- 0,\y);

    % Realtime notification
    \def\y{-11.4}
    \draw[msg, blue!60] (pg |- 0,\y) -- node[above, blue!60] {Realtime: route\_assigned} (driver |- 0,\y);
    \def\y{-12.2}
    \draw[ret] (flutter |- 0,\y) -- node[above] {Map preview + driver rec.} (admin |- 0,\y);
  \end{tikzpicture}}%
  \caption{Sequence diagram: task creation via CSV import, geocoding,
           route generation, and real-time driver notification.}
  \label{fig:seq-task}
\end{figure}

\subsection{Agentic AI Suggestion Cycle}
\label{sec:seq-agentic}

Figure~\ref{fig:seq-agentic} details the interaction between the driver
application, the AI service, and the remote LLM when producing an agentic
suggestion with a one-tap executable action.

\begin{figure}[ht]
  \centering
  \resizebox{0.92\textwidth}{!}{%
  \begin{tikzpicture}[
    actor/.style={draw, rectangle, rounded corners=3pt, minimum width=2cm,
                  minimum height=0.7cm, font=\sffamily\footnotesize,
                  fill=#1!15, align=center, thick},
    msg/.style={-{Stealth[length=2mm]}, thick, font=\sffamily\tiny},
    ret/.style={-{Stealth[length=2mm]}, thick, dashed, font=\sffamily\tiny},
    note/.style={draw, rectangle, rounded corners=2pt, fill=yellow!10,
                 font=\sffamily\tiny, align=left, inner sep=3pt},
    alt/.style={draw, dashed, rounded corners=2pt, fill=green!5, inner sep=4pt},
  ]
    % Participants
    \node[actor=teal]   (driver)  at (0,0)    {Driver\\UI};
    \node[actor=orange]  (taskvm)  at (3.5,0)  {Task\\ViewModel};
    \node[actor=violet]  (aisvc)   at (7,0)    {AI\\Service};
    \node[actor=blue]   (openr)   at (10.5,0) {OpenRouter\\API};
    \node[actor=blue]   (supa)    at (14,0)   {Supabase\\(state)};

    % Lifelines
    \foreach \n in {driver, taskvm, aisvc, openr, supa} {
      \draw[dashed, gray!40] (\n.south) -- ++(0,-14);
    }

    % Messages
    \def\y{-1.2}
    \draw[msg] (driver |- 0,\y) -- node[above] {Tap ``AI Suggest''} (taskvm |- 0,\y);
    \def\y{-2.0}
    \draw[msg] (taskvm |- 0,\y) -- node[above] {getAgenticSuggestion(context)} (aisvc |- 0,\y);
    \def\y{-2.8}
    \node[note] at (7, -3.3) {System prompt +\\JSON schema +\\task context};
    \draw[msg] (aisvc |- 0,\y) -- node[above] {POST /chat/completions} (openr |- 0,\y);
    \def\y{-4.0}
    \draw[ret] (openr |- 0,\y) -- node[above] {\{suggestion, action, taskId\}} (aisvc |- 0,\y);

    % Parse
    \def\y{-5.0}
    \draw[msg] (aisvc |- 0,\y) -- node[above] {parseJson $\to$ AiAction} (aisvc |- 0,\y);
    \def\y{-5.8}
    \draw[ret] (aisvc |- 0,\y) -- node[above] {AiAction object} (taskvm |- 0,\y);
    \def\y{-6.6}
    \draw[ret] (taskvm |- 0,\y) -- node[above] {Render chip: suggestion + action} (driver |- 0,\y);

    % User decision
    \def\y{-7.8}
    \node[note] at (-1.8, \y) {Human-in-the-loop\\confirmation gate};
    \draw[msg] (driver |- 0,\y) -- node[above] {Tap ``Confirm''} (taskvm |- 0,\y);
    \def\y{-8.6}
    \draw[msg] (taskvm |- 0,\y) -- node[above] {executeAction(action)} (taskvm |- 0,\y);

    % State update
    \def\y{-9.4}
    \draw[msg] (taskvm |- 0,\y) -- node[above] {UPDATE task status} (supa |- 0,\y);
    \def\y{-10.2}
    \draw[ret] (supa |- 0,\y) -- node[above] {OK (RLS verified)} (taskvm |- 0,\y);
    \def\y{-11.0}
    \draw[ret] (taskvm |- 0,\y) -- node[above] {UI refresh (new state)} (driver |- 0,\y);

    % Alt: reject
    \def\y{-12.2}
    \node[note] at (-1.8, \y) {Alt: user taps\\``Dismiss'' $\to$ no-op};
  \end{tikzpicture}}%
  \caption{Sequence diagram: agentic AI suggestion with human-in-the-loop
           confirmation and state synchronisation.}
  \label{fig:seq-agentic}
\end{figure}

The key design insight is the \emph{confirmation gate} between steps~6 and~7:
the LLM's suggestion is never executed automatically.  This addresses the
inherent unreliability of LLM outputs in safety-critical logistics operations
\autocite{ji2024hallucination} and distinguishes OrderVista's agentic approach
from fully autonomous AI agents.
